\section{Scope, principal results, and derivational order}\label{sec:intro}
\subsection{The mathematical problem}
A deterministic law specifies the future of a complete state under a fixed intervention. It does not, by itself, specify a laboratory preparation, a measurement resolution, a family of available controls, or a rule for evaluating unresolved outcomes. We therefore distinguish four primitives: mechanical evolution, preparation, intervention, and observation. A path distribution is constructed from these primitives. A nonlinear evaluation is then attached to a specified selection or decision rule. A reduced state becomes legitimate only after its predictive and operational sufficiency has been proved.

The objective is not a universal diffusion theorem for every mechanical system. It is a framework in which a proposed mechanical-to-statistical implication can be checked without inserting its desired conclusion into the setup. The basic path construction requires no independent increments, mixing rate, spectral gap, or Markov property of the observed coordinates. The response theorems impose smoothness and integrability exactly where differentiation is used. The continuous-alphabet prediction theorem has stronger, separately stated hypotheses than the measurable path construction.

The term \emph{realizable} always refers to declared apparatus and preparation operations. A formula for a density is not an implementation. Conversely, a finite-energy force may change the support of a deterministic path law, so its physical cost is not automatically a relative entropy. We retain both distinctions throughout.

\subsection{Principal statements}
The main results form six groups, each proved below.

\paragraph{Realizability and coarse preparation.}
Theorem~\ref{thm:real-gap} identifies the exact loss incurred by restricting a Gibbs variational problem to attainable path laws. Theorem~\ref{thm:coarse-prep} solves the restriction in which only a specified initial statistic can be reweighted. Theorem~\ref{thm:cap} solves a bounded-density version. Proposition~\ref{prop:acceptance} states precisely which postselection and preparation operations implement these laws and what independent-trial budget they consume. No information cost is identified with mechanical work.

\paragraph{Joint instruments.}
Theorem~\ref{thm:instrument-path} constructs chronological records from a joint state--observation instrument. Proposition~\ref{prop:pointer} gives an explicit Hamiltonian pointer shear with both readout and back-action. The apparatus is part of the complete mechanical state; repeated measurements do not become independent unless their preparation and coupling justify that property.

\paragraph{Prediction and resources.}
Theorems~\ref{thm:finite-state} and \ref{thm:response-state} construct a complete finite-alphabet predictor and its parameter-response extension. Theorem~\ref{thm:continuous-state} proves a compact continuous-alphabet version from jointly continuous, positive finite-record densities. Theorems~\ref{thm:budget-dpp} and \ref{thm:info-dpp} prove dynamic programming only after residual resources and legal pasting have been included.

\paragraph{Observable-dependent errors.}
Theorems~\ref{thm:test-tilt} and \ref{thm:jet-error} control normalized, tilted, and differentiated experimental quantities by their actual unnormalized tests. These estimates do not require total-variation convergence of singular microscopic path laws. Theorem~\ref{thm:value-error} gives a separate finite-horizon value bound and specifies the continuation tests that have to be controlled.

\paragraph{Geometric response.}
Theorems~\ref{thm:shape} and \ref{thm:chamber} combine moving-domain derivatives, trajectory jets, collision-time derivatives, and saltation. Theorem~\ref{thm:assembly} provides a complete-preparation Cauchy criterion rather than inferring derivative convergence from a small exceptional probability.

\paragraph{A whole-preparation mechanical calculation.}
Theorem~\ref{thm:global-tube} proves an exact collision-tube identity for the entire natural equilibrium preparation of a periodic triangular Lorentz gas over a short physical horizon. Theorem~\ref{thm:global-response} differentiates the resulting physical path distributions through every fixed finite order. Arbitrarily grazing incidences are included in the integrals. The construction works because event-time coordinates produce a flux factor rather than an inverse normal velocity, and because a strict intercollision gap excludes a second collision within this horizon. This is stronger than a prepared regular trajectory tube, but it is not a theorem about arbitrary collision histories at long times.

\subsection{Classical inputs and proposed contribution}
We use standard Borel disintegration and kernel extension as foundational measure theory \cite{Kallenberg}. The Gibbs formula, entropy chain rule, and positive exponential character are proved explicitly, but are not presented as newly discovered probability laws. Characterizations of relative entropy in \cite{BaezFritz,GagnePanangaden} have their own information-composition assumptions; we do not claim that mechanics alone selects relative entropy.

Predictive representations of state have a substantial history \cite{LittmanSuttonSingh}. Our all-tests representation is not a finite-rank realization or a learning algorithm. General partially observed control also requires care with posterior continuity \cite{FeinbergEtAl}; our continuous prediction result states sufficient density conditions and proves the quotient and update rather than importing those properties from a finite-alphabet formula.

Saltation is a standard hybrid sensitivity tool \cite{KongEtAl}. The mechanical use here is in its combination with terminal-interface currents and normalized experimental quantities. Perturbation theory for the Lorentz gas \cite{DemersZhang} is not treated as a blanket source of arbitrary-order moving-table response. Likewise, driven-process theory \cite{ChetriteTouchette,ChetriteTouchetteControl} does not identify every finite bridge with an available force protocol. The role of legal pasting in our resource construction is related to dynamic-consistency questions \cite{EpsteinSchneider}, but the finite-horizon claims below have direct proofs.

The proposed research contribution is the explicit compatibility of realizability restrictions, information loss, resource-aware prediction, and geometric response, together with the full-preparation finite-time collision assembly. Priority and journal significance require independent evaluation. Computation is used to check finite identities and source integrity, not to certify an infinite family of theorems.

\subsection{Notation and quantifiers}
All unspecified measurable spaces are standard Borel. Probability measures on $E$ form $\cP(E)$. For a finite signed measure $\lambda$,
\[
 \norm\lambda_{\rm var}=\sup_{|f|\le1}\left|\int f\dd\lambda\right|,
 \qquad d_{\TV}(P,Q)=\tfrac12\norm{P-Q}_{\rm var}.
\]
Thus $|Pf-Qf|\le\osc(f)d_{\TV}(P,Q)$. Relative entropy is
$D(Q\Vert P)=\int q\log q\dd P$ for $Q=qP$, and $+\infty$ otherwise; $0\log0=0$. Physical parameters are denoted $a$, or $R$ for a radius. A source is $\xi$ or $V$, and an evaluation intensity is $\eta$. Equality of these objects must be justified, not inferred from similar notation.

An equality of conditional laws is initially an almost-everywhere statement for its conditioning law. Continuous versions, uniformity in parameters, and comparisons on null sets are additional properties established when needed. A finite-order statement means every prescribed finite $r$ under the corresponding regularity; it does not imply one analytic infinite-order bound. Whole-preparation response below refers to the specified equilibrium preparation and horizon, not to every preparation and time scale.
\par
